\documentclass[12pt]{article}
\usepackage[margin=0.6in,top=0.85in,bottom=0.55in]{geometry}
\usepackage{lmodern}
\usepackage{microtype}
\usepackage{fancyhdr}
\usepackage{titlesec}
\usepackage{enumitem}
\usepackage{lastpage}
\usepackage[hidelinks]{hyperref}

\setlength{\parindent}{0pt}
\setlength{\parskip}{0.25em}
\setlength{\headheight}{26pt}
\titleformat{\section}{\large\bfseries\scshape}{}{0pt}{}
\pagestyle{fancy}
\fancyhf{}
\fancyhead[C]{\footnotesize Student Name: \hspace{2.3in} \quad Assignment ID: U1L9LE \quad Page \thepage\ of \pageref{LastPage}}
\renewcommand{\headrulewidth}{0.5pt}

\newcommand{\answerbox}[2]{%
\vspace{0.12em}
\noindent\textbf{#1}\par
\vspace{0.04em}
\noindent\fbox{%
\begin{minipage}[t][#2]{\dimexpr\textwidth-2\fboxsep-2\fboxrule}
\rule{\textwidth}{0pt}
\vspace{#2}
\end{minipage}}
\par\vspace{0.22em}}

\begin{document}

\begin{center}
{\LARGE\bfseries Week 9 Lit Example Worksheet}\par\vspace{0.05em}
{\large\scshape Julius Caesar: When the Reason Assumes the Verdict}\par\vspace{0.12em}
\rule{0.78\textwidth}{0.5pt}
\end{center}

\textbf{Purpose:} Apply Week 9's method to Brutus's reasoning. Find where the conclusion is assumed and rewrite for independent support.

\textbf{Directions:} Use the Lit Example Reader. Do not collapse circularity with mere emotional force.

\section*{Part 1: Brutus Alone}

\textbf{Part 1 Q1:} What conclusion does Brutus announce at the opening of Passage 1?

\answerbox{Part 1 Q1 ANSWER BOX}{2.0in}

\textbf{Part 1 Q2:} What present evidence about Caesar does Brutus admit is weak or missing, and what future danger does he rely on instead?

\answerbox{Part 1 Q2 ANSWER BOX}{2.45in}

\newpage

\section*{Part 2: Circularity Test}

\textbf{Part 2 Q1:} Could a fair opponent accept Brutus's main reasons without already accepting that Caesar must die? Explain.

\answerbox{Part 2 Q1 ANSWER BOX}{2.45in}

\textbf{Part 2 Q2:} Where does ``Fashion it thus'' or the serpent's-egg comparison risk assuming the verdict?

\answerbox{Part 2 Q2 ANSWER BOX}{2.35in}

\newpage

\section*{Part 3: Public Speech}

\textbf{Part 3 Q1:} In Passage 2, how can the ambition $\rightarrow$ death move become circular or question-begging if ambition is not independently proved?

\answerbox{Part 3 Q1 ANSWER BOX}{2.45in}

\textbf{Part 3 Q2:} How does the slaves/free-men alternative pressure the crowd in a way that may assume the danger still needing proof?

\answerbox{Part 3 Q2 ANSWER BOX}{2.45in}

\newpage

\section*{Part 4: Rewrite}

\textbf{Part 4 Q1:} Rewrite Brutus's case in non-circular form: keep a clear claim and state what independent evidence would be required before death is justified.

\answerbox{Part 4 Q1 ANSWER BOX}{2.55in}

\textbf{Part 4 Q2:} In one paragraph, explain how Week 9's tool helps a reader separate a coherent speech from a proved conclusion in these passages.

\answerbox{Part 4 Q2 ANSWER BOX}{2.45in}

\end{document}
